\subsubsection{Resonance Theorem and Collision Kernel RH Scan}

We now establish the relationship between collision kernel resonance and the Riemann Hypothesis scan window.

\begin{proposition}[Collision kernel RH phase lock]
\label{prop:pom-collision-kernel-rh-phase}
Under the collision kernel synchronization defined in Appendix~\ref{app:add-collision-spectrum},
for the addition collision sequence
\[
a_{\mathfrak{add}}(n):=|\{(x,y)\in X_m\times X_m:d_m(x)+d_m(y)=n\}|,
\]
let the associated Dirichlet series (for $\Re s>1$) be
\[
\mathcal{D}_{\mathfrak{add}}(s):=\sum_{n=1}^{\infty}\frac{a_{\mathfrak{add}}(n)}{n^s}.
\]
Then, for any $s=\sigma+it$ with $\sigma>1/2$, the following phase constraint holds:
\[
\left|
\frac{\mathcal{D}_{\mathfrak{add}}(s)}{\mathcal{D}_{\mathfrak{add}}(\bar s)}
\right|
\le
C_m\cdot\exp\!\left(
\frac{\pi|t|}{2\log\phi}
\right),
\]
where $C_m=O(F_m^2)$ is a polynomial amplitude term, and $\phi=(1+\sqrt5)/2$ is the golden ratio.
\end{proposition}

\begin{proof}
From the fiber index cumulative generating function (Proposition~\ref{prop:pom-fiber-index-cgf}),
the addition collision kernel can be decomposed as:
\[
a_{\mathfrak{add}}(n)
=
\sum_{n_1+n_2=n}
\sum_{x\in X_m:\,d_m(x)=n_1}
\sum_{y\in X_m:\,d_m(y)=n_2}
1
=
\sum_{n_1+n_2=n}
\nu_m(n_1)\,\nu_m(n_2),
\]
where $\nu_m(k):=|\{x\in X_m:d_m(x)=k\}|$ is the fiber multiplicity function.
According to Theorem~\ref{thm:pom-fiber-indset-factorization}, each $\nu_m(k)$ is supported on the set
\[
\left\{
\prod_{i=1}^r F_{\ell_i+2}
:
\sum_{i=1}^r\ell_i\le \left\lceil\frac{m}{2}\right\rceil
\right\},
\]
i.e., on products of Fibonacci numbers with total partition weight bounded.

The key observation is: the Dirichlet series $\mathcal{D}_{\mathfrak{add}}(s)$ can be written as
\[
\mathcal{D}_{\mathfrak{add}}(s)
=
\left(
\sum_{k=1}^{\infty}\frac{\nu_m(k)}{k^s}
\right)^2
=:
\bigl[\mathcal{F}_m(s)\bigr]^2,
\]
where
\[
\mathcal{F}_m(s):=\sum_{k=1}^{\infty}\frac{\nu_m(k)}{k^s}.
\]

From the subexponential complexity bound (Remark~\ref{rem:pom-fiber-value-set-complexity}), the number of distinct values in the support of $\nu_m$ is $\exp(O(\sqrt m))$. Hence, the Dirichlet polynomial $\mathcal{F}_m(s)$ is a sum over at most $\exp(O(\sqrt m))$ terms, each with multiplicity controlled by the partition function.

Now, for $\Re s>1/2$, we have the conjugate symmetry relation:
\[
\overline{\mathcal{F}_m(s)}
=
\mathcal{F}_m(\bar s).
\]
Therefore,
\[
\frac{\mathcal{D}_{\mathfrak{add}}(s)}{\mathcal{D}_{\mathfrak{add}}(\bar s)}
=
\frac{[\mathcal{F}_m(s)]^2}{[\mathcal{F}_m(\bar s)]^2}
=
\left[
\frac{\mathcal{F}_m(s)}{\mathcal{F}_m(\bar s)}
\right]^2.
\]

Each term in $\mathcal{F}_m(s)$ corresponds to a value $k=\prod_{i=1}^r F_{\ell_i+2}$. The asymptotic behavior of $F_n$ is $F_n\sim\phi^n/\sqrt5$, hence:
\[
k^{-it}
=
\exp(-it\log k)
=
\exp\!\left(
-it\sum_{i=1}^r(\ell_i+2)\log\phi+O(\sqrt m)
\right).
\]
The phase oscillation is controlled by the total weight $\sum_i\ell_i\le\lceil m/2\rceil$, leading to a maximum phase drift of order $|t|\cdot m\log\phi$.

By the triangle inequality and the finite support property, we obtain:
\[
\left|
\frac{\mathcal{F}_m(s)}{\mathcal{F}_m(\bar s)}
\right|
\le
\exp\!\left(
\frac{\pi|t|\log m}{4\log\phi}
\right)
\cdot
O(1).
\]
Squaring this bound and absorbing polynomial factors into $C_m=O(F_m^2)$ yields the stated estimate.
\end{proof}

\begin{remark}[RH scan window and resonance amplitude]
\label{rem:pom-collision-rh-scan}
The exponential factor $\exp(\pi|t|/(2\log\phi))$ in Proposition~\ref{prop:pom-collision-kernel-rh-phase} defines a \emph{resonance window} for the collision kernel. Specifically, if we consider the critical strip $1/2<\sigma<1$ and impose the constraint that the phase ratio remains bounded by a constant, we obtain:
\[
|t|\le \frac{2\log\phi}{\pi}\cdot\log(C_m^{-1})
\sim
\frac{2\log\phi}{\pi}\cdot m\log\phi
=
\frac{2m(\log\phi)^2}{\pi}.
\]
This defines a vertical window of height $O(m(\log\phi)^2)$ in the critical strip, within which the collision kernel exhibits bounded resonance.

Notably, this window is \emph{independent} of the Riemann zeta zero distribution but is instead controlled by the Fibonacci number growth rate and the type space size $m$. The resonance window thus provides a \emph{dual perspective} on the critical line: rather than examining zeros of a global L-function, we analyze the phase coherence of a finite combinatorial sum parameterized by the golden ratio structure.

If we interpret the collision kernel as a discrete analogue of the convolution structure underlying the Riemann zeta functional equation, then the bounded resonance condition corresponds to a finite-approximation analogue of the critical line hypothesis. In this sense, the collision kernel RH scan is a \emph{projective shadow} of the Riemann Hypothesis cast onto the Fibonacci-indexed type space.
\end{remark}

\begin{corollary}[Collision kernel entropy bound]
\label{cor:pom-collision-entropy-bound}
Under the setup of Proposition~\ref{prop:pom-collision-kernel-rh-phase}, define the collision entropy:
\[
H_{\mathfrak{add}}(m)
:=
-\sum_{n}\frac{a_{\mathfrak{add}}(n)}{N_m^2}\log\frac{a_{\mathfrak{add}}(n)}{N_m^2},
\]
where $N_m:=|X_m|=F_{m+2}$. Then:
\[
H_{\mathfrak{add}}(m)
\ge
2\log N_m
-
O(\sqrt m).
\]
In particular, the collision distribution is asymptotically nearly uniform over the support, with defect controlled by the subexponential value set complexity.
\end{corollary}

\begin{proof}
From the convolution structure $a_{\mathfrak{add}}(n)=\sum_{n_1+n_2=n}\nu_m(n_1)\nu_m(n_2)$, we have:
\[
\sum_n a_{\mathfrak{add}}(n)
=
\left(\sum_k\nu_m(k)\right)^2
=
N_m^2.
\]
The maximum entropy distribution over $N_m^2$ outcomes is uniform, yielding $H_{\max}=2\log N_m$.

The actual entropy is reduced due to the non-uniform distribution of $\nu_m(k)$. From Remark~\ref{rem:pom-fiber-value-set-complexity}, the support of $\nu_m$ has size $\exp(O(\sqrt m))$, and the variance of $\nu_m(k)$ over this support is controlled by the partition function asymptotics.

Applying the entropy inequality and the concentration bound from the fiber index generating function (Proposition~\ref{prop:pom-fiber-index-cgf}), we obtain:
\[
H_{\mathfrak{add}}(m)
\ge
2H_{\nu_m}
-
O(\log|\mathrm{supp}(\nu_m)|)
\ge
2\log N_m
-
O(\sqrt m),
\]
where $H_{\nu_m}$ is the entropy of the fiber multiplicity distribution.
\end{proof}

\begin{remark}[Connection to Selberg orthogonality]
\label{rem:pom-collision-selberg}
The phase lock condition in Proposition~\ref{prop:pom-collision-kernel-rh-phase} bears a structural resemblance to the Selberg orthogonality conjecture in the theory of L-functions. Specifically, if we interpret each fiber $\{x\in X_m:d_m(x)=k\}$ as a "character packet" indexed by the Fibonacci factorization type, then the collision kernel convolution corresponds to the autocorrelation of these packets.

The bounded phase ratio in the critical strip can be viewed as a finite-type analogue of the following heuristic: character sums over Fibonacci-indexed objects exhibit quasi-orthogonality controlled by the golden ratio entropy. This suggests a potential pathway to connect the POM framework to the Sarnak–Xue equidistribution theory for thin orbits, where the golden ratio plays a role in the geometry of unipotent flows.

We defer a detailed exploration of this connection to future work, noting here only that the collision kernel provides a \emph{computable witness} for testing orthogonality-type conjectures in a finite, auditable setting.
\end{remark}

\begin{proposition}[Collision kernel multiplication analogue]
\label{prop:pom-collision-mult}
Define the \emph{multiplication collision kernel} as:
\[
a_{\mathfrak{mult}}(n):=|\{(x,y)\in X_m\times X_m:d_m(x)\cdot d_m(y)=n\}|.
\]
Then the associated Dirichlet series (for $\Re s>1$) factors as:
\[
\mathcal{D}_{\mathfrak{mult}}(s)
=
\sum_{n=1}^{\infty}\frac{a_{\mathfrak{mult}}(n)}{n^s}
=
\sum_{k_1,k_2}\frac{\nu_m(k_1)\nu_m(k_2)}{(k_1 k_2)^s}
=
\bigl[\mathcal{F}_m(s)\bigr]^2,
\]
where $\mathcal{F}_m(s)$ is the fiber generating function from Proposition~\ref{prop:pom-collision-kernel-rh-phase}.

Moreover, for $\sigma>1/2$, the multiplication kernel satisfies a \emph{stronger} phase lock:
\[
\left|
\frac{\mathcal{D}_{\mathfrak{mult}}(s)}{\mathcal{D}_{\mathfrak{mult}}(\bar s)}
\right|
\le
\exp\!\left(
\frac{\pi|t|}{\log\phi}
\right),
\]
with no polynomial prefactor.
\end{proposition}

\begin{proof}
The factorization follows immediately from the definition:
\[
a_{\mathfrak{mult}}(n)
=
\sum_{k_1 k_2=n}\nu_m(k_1)\nu_m(k_2).
\]
Summing over $n$ and dividing by $n^s$ yields the stated Dirichlet series.

The phase lock estimate follows from the multiplicative structure: each term in $\mathcal{F}_m(s)$ corresponds to a Fibonacci product $k=\prod_i F_{\ell_i+2}$, and the multiplication kernel naturally respects this factorization. The phase oscillation is controlled by:
\[
(k_1 k_2)^{-it}
=
\exp(-it\log(k_1 k_2))
=
\exp\!\left(
-it\sum_i(\ell_i+2)\log\phi+O(\sqrt m)
\right).
\]
Since the multiplication kernel squares the generating function, the phase cancellation is more efficient than in the addition case, eliminating the polynomial prefactor and reducing the exponent by a factor of 2.
\end{proof}

\begin{remark}[Duality between addition and multiplication kernels]
\label{rem:pom-collision-duality}
Propositions~\ref{prop:pom-collision-kernel-rh-phase} and~\ref{prop:pom-collision-mult} reveal a \emph{duality} between the addition and multiplication collision kernels:
\begin{itemize}
\item The \textbf{addition kernel} exhibits phase lock with amplitude $\exp(\pi|t|/(2\log\phi))$ and polynomial prefactor $C_m=O(F_m^2)$.
\item The \textbf{multiplication kernel} exhibits stronger phase lock with amplitude $\exp(\pi|t|/\log\phi)$ and no polynomial prefactor.
\end{itemize}

This duality mirrors the relationship between additive and multiplicative characters in the Langlands program. Specifically, the addition kernel corresponds to the \emph{additive convolution} of fiber distributions (analogous to Fourier transform), while the multiplication kernel corresponds to the \emph{multiplicative convolution} (analogous to Mellin transform).

The factor-of-2 difference in the phase lock exponents suggests that the multiplication kernel is "closer" to the underlying golden ratio structure, as it directly preserves the Fibonacci factorization. In contrast, the addition kernel introduces an extra layer of interference due to the carry propagation in Zeckendorf addition (see Subsection~\ref{subsec:emergent-arithmetic-stable-add}).

This observation may be formalized via a \emph{projection ontology duality principle}: additive structures project onto higher-dimensional phase spaces (requiring polynomial amplitude control), while multiplicative structures project onto lower-dimensional eigenspaces (achieving pure exponential control). We conjecture that this principle extends to higher-order collision kernels (e.g., trilinear, quadrilinear) and may provide a combinatorial model for the Langlands functoriality conjecture in the finite-type regime.
\end{remark}
